% questions.tex —— 2013 年·广东卷·文科 · 题目内容片段（无答案）
% 由 tools/generate_projects.py 依来源总账自 2013/2013_数学真题_广东文_试题.md 生成，勿手改（改源 MD 后重跑）

\gkpart{一、选择题}

\begin{question}
\gkti{1}设集合 $S=\{x\mid x^2+2x=0,\ x\in\mathbf R\}$，$T=\{x\mid x^2-2x=0,\ x\in\mathbf R\}$，则 $S\cap T=$\mbox{（　　）}

（A）$\{0\}$；（B）$\{0,2\}$；（C）$\{-2,0\}$；（D）$\{-2,0,2\}$。

\end{question}
\begin{question}
\gkti{2}函数 $y=\dfrac{\lg(x+1)}{x-1}$ 的定义域是\mbox{（　　）}

（A）$(-1,+\infty)$；（B）$[-1,+\infty)$；（C）$(-1,1)\cup(1,+\infty)$；（D）$[-1,1)\cup(1,+\infty)$。

\end{question}
\begin{question}
\gkti{3}若 $\mathrm i(x+y\mathrm i)=3+4\mathrm i$，$x$，$y\in\mathbf R$，则复数 $x+y\mathrm i$ 的模是\mbox{（　　）}

（A）$2$；（B）$3$；（C）$4$；（D）$5$。

\end{question}
\begin{question}
\gkti{4}已知 $\sin\left(\dfrac{5\pi}2+\alpha\right)=\dfrac15$，那么 $\cos\alpha=$\mbox{（　　）}

（A）$-\dfrac25$；（B）$-\dfrac15$；（C）$\dfrac15$；（D）$\dfrac25$。

\Needspace{0.4\textheight}
\end{question}
\begin{question}
\gkti{5}执行如图所示的程序框图，若输入 $n$ 的值为 $3$，则输出 $s$ 的值是\mbox{（　　）}
\nopagebreak[4]
\begin{center}% 2013 广东卷文科第 5 题：程序框图（输入 n，累加 s=s+(i-1)）。
\begin{tikzpicture}[font=\footnotesize, >=Stealth, line width=0.4pt,
  node distance=3.2mm,
  box/.style={draw, minimum height=1.65em, inner xsep=6pt},
  term/.style={draw, rounded corners=6pt, minimum height=1.65em, inner xsep=9pt},
  io/.style={draw, trapezium, trapezium left angle=70, trapezium right angle=110,
             minimum height=1.65em, inner xsep=4pt},
  dec/.style={draw, diamond, aspect=2.5, inner sep=1.2pt}]
  \node[term] (start) {开始};
  \node[io, below=5mm of start] (input) {输入 $n$};
  \node[box, below=5mm of input] (init) {$i=1$，$s=1$};
  \node[dec, below=6mm of init] (cond) {$i\leqslant n$};
  \node[box, below=8mm of cond] (acc) {$s=s+(i-1)$};
  \node[box, below=5mm of acc] (inc) {$i=i+1$};
  \node[io, right=13mm of cond] (output) {输出 $s$};
  \node[term, below=6mm of output] (end) {结束};
  \draw[->] (start) -- (input);
  \draw[->] (input) -- (init);
  \draw[->] (init) -- (cond);
  \draw[->] (cond) -- node[right] {是} (acc);
  \draw[->] (acc) -- (inc);
  \draw[->] (inc.west) -- ++(-1.15,0) |- (cond.west);
  \draw[->] (cond.east) -- node[above] {否} (output.west);
  \draw[->] (output) -- (end);
\end{tikzpicture}
\end{center}
\nopagebreak[4]
（A）$1$；（B）$2$；（C）$4$；（D）$7$。

\Needspace{0.4\textheight}
\end{question}
\begin{question}
\gkti{6}某三棱锥的三视图如图所示，则该三棱锥的体积是\mbox{（　　）}
\nopagebreak[4]
\begin{center}% 2013 广东卷文科第 6 题：三棱锥的三视图（正视图、侧视图、俯视图）。
\begin{tikzpicture}[font=\scriptsize, >=Stealth, line width=0.4pt,
  x=0.95cm, y=0.95cm]
  % 正视图：直角三角形，竖直直角边在左，高 2、底 1
  \begin{scope}
    \draw (0,0) -- (1,0) -- (0,2) -- cycle;
    \draw[<->] (-0.42,0) -- (-0.42,2) node[midway,left] {$2$};
    \draw[<->] (0,-0.38) -- (1,-0.38) node[midway,below] {$1$};
    \node at (0.5,-1.15) {正视图};
  \end{scope}
  % 侧视图：直角三角形，竖直直角边在右，高 2、底 1
  \begin{scope}[xshift=3.05cm]
    \draw (0,0) -- (1,0) -- (1,2) -- cycle;
    \draw[<->] (0,-0.38) -- (1,-0.38) node[midway,below] {$1$};
    \node at (0.5,-1.15) {侧视图};
  \end{scope}
  % 俯视图：等腰直角三角形，两条直角边均为 1
  \begin{scope}[yshift=-3.55cm]
    \draw (0,0) -- (1,0) -- (1,1) -- cycle;
    \node at (0.5,-1.05) {俯视图};
  \end{scope}
\end{tikzpicture}
\end{center}
\nopagebreak[4]
（A）$\dfrac16$；（B）$\dfrac13$；（C）$\dfrac23$；（D）$1$。

\end{question}
\begin{question}
\gkti{7}垂直于直线 $y=x+1$ 且与圆 $x^2+y^2=1$ 相切于第一象限的直线方程是\mbox{（　　）}

（A）$x+y-\sqrt2=0$；（B）$x+y+1=0$；（C）$x+y-1=0$；（D）$x+y+\sqrt2=0$。

\end{question}
\begin{question}
\gkti{8}设 $l$ 为直线，$\alpha$，$\beta$ 是两个不同的平面，下列命题中正确的是\mbox{（　　）}

（A）若 $l\parallel\alpha$，$l\parallel\beta$，则 $\alpha\parallel\beta$；（B）若 $l\perp\alpha$，$l\perp\beta$，则 $\alpha\parallel\beta$；（C）若 $l\perp\alpha$，$l\parallel\beta$，则 $\alpha\parallel\beta$；（D）若 $\alpha\perp\beta$，$l\parallel\alpha$，则 $l\perp\beta$。

\end{question}
\begin{question}
\gkti{9}已知中心在原点的椭圆 $C$ 的右焦点为 $F(1,0)$，离心率等于 $\dfrac12$，则 $C$ 的方程是\mbox{（　　）}

（A）$\dfrac{x^2}3+\dfrac{y^2}4=1$；（B）$\dfrac{x^2}4+\dfrac{y^2}{\sqrt3}=1$；（C）$\dfrac{x^2}4+\dfrac{y^2}2=1$；（D）$\dfrac{x^2}4+\dfrac{y^2}3=1$。

\end{question}
\begin{question}
\gkti{10}设 $\boldsymbol a$ 是已知的平面向量且 $\boldsymbol a\neq\boldsymbol 0$。关于向量 $\boldsymbol a$ 的分解，有如下四个命题：

① 给定向量 $\boldsymbol b$，总存在向量 $\boldsymbol c$，使 $\boldsymbol a=\boldsymbol b+\boldsymbol c$；

② 给定向量 $\boldsymbol b$ 和 $\boldsymbol c$，总存在实数 $\lambda$ 和 $\mu$，使 $\boldsymbol a=\lambda\boldsymbol b+\mu\boldsymbol c$；

③ 给定单位向量 $\boldsymbol b$ 和正数 $\mu$，总存在单位向量 $\boldsymbol c$ 和实数 $\lambda$，使 $\boldsymbol a=\lambda\boldsymbol b+\mu\boldsymbol c$；

④ 给定正数 $\lambda$ 和 $\mu$，总存在单位向量 $\boldsymbol b$ 和单位向量 $\boldsymbol c$，使 $\boldsymbol a=\lambda\boldsymbol b+\mu\boldsymbol c$。

上述命题中的向量 $\boldsymbol b$，$\boldsymbol c$ 和 $\boldsymbol a$ 在同一平面内且两两不共线，则真命题的个数是\mbox{（　　）}

（A）$1$；（B）$2$；（C）$3$；（D）$4$。

\end{question}
\gkpart{二、填空题}

\begin{question}
\gkti{11}设数列 $\{a_n\}$ 是首项为 $1$，公比为 $-2$ 的等比数列，则 $a_1+|a_2|+a_3+|a_4|=$\gkblank{}。

\end{question}
\begin{question}
\gkti{12}若曲线 $y=ax^2-\ln x$ 在点 $(1,a)$ 处的切线平行于 $x$ 轴，则 $a=$\gkblank{}。

\end{question}
\begin{question}
\gkti{13}已知变量 $x$，$y$ 满足约束条件 $\begin{cases}x-y+3\geqslant0,\\-1\leqslant x\leqslant1,\\y\geqslant1,\end{cases}$ 则 $z=x+y$ 的最大值是\gkblank{}。

\end{question}
\begin{question}
\gkti{14}已知曲线 $C$ 的极坐标方程为 $\rho=2\cos\theta$，以极点为原点，极轴为 $x$ 轴的正半轴建立直角坐标系，则曲线 $C$ 的参数方程为\gkblank{}。

\Needspace{0.4\textheight}
\end{question}
\begin{question}
\gkti{15}如图，在矩形 $ABCD$ 中，$AB=\sqrt3$，$BC=3$，$BE\perp AC$，垂足为 $E$，则 $ED=$\gkblank{}。
\nopagebreak[4]
\begin{center}% 2013 广东卷文科第 15 题：矩形 ABCD，AB=sqrt3，BC=3，BE 垂直 AC 于 E。
\begin{tikzpicture}[font=\small, >=Stealth, line width=0.4pt]
  \coordinate (A) at (0,0);
  \coordinate (B) at (0,2.2517);
  \coordinate (C) at (3.9,2.2517);
  \coordinate (D) at (3.9,0);
  \coordinate (E) at (0.975,0.5629);
  \draw (A) -- (B) -- (C) -- (D) -- cycle;
  \draw (A) -- (C);
  \draw (B) -- (E);
  \draw (E) -- (D);
  \node[below left] at (A) {$A$};
  \node[above left] at (B) {$B$};
  \node[above right] at (C) {$C$};
  \node[below right] at (D) {$D$};
  \node[above right] at (E) {$E$};
\end{tikzpicture}
\end{center}
\nopagebreak[4]
\end{question}
\gkpart{三、解答题}

\begin{question}
\gkti{16}已知函数 $f(x)=\sqrt2\cos\left(x-\dfrac{\pi}{12}\right)$，$x\in\mathbf R$。

（1） 求 $f\left(\dfrac{\pi}3\right)$ 的值；

（2） 若 $\cos\theta=\dfrac35$，$\theta\in\left(\dfrac{3\pi}2,2\pi\right)$，求 $f\left(\theta-\dfrac{\pi}6\right)$。

\end{question}
\begin{question}
\gkti{17}从一批苹果中，随机抽取 $50$ 个，其重量（单位：克）的频数分布表如下：

\begin{center}
\begin{tabular}{lllll}
\toprule
分组（重量） & $[80,85)$ & $[85,90)$ & $[90,95)$ & $[95,100)$ \\
\midrule
频数 & $5$ & $10$ & $20$ & $15$ \\
\bottomrule
\end{tabular}
\end{center}

（1） 根据频数分布表计算苹果的重量在 $[90,95)$ 的频率；

（2） 用分层抽样的方法从重量在 $[80,85)$ 和 $[95,100)$ 的苹果中共抽取 $4$ 个，其中重量在 $[80,85)$ 的有几个？

（3） 在（2）中抽出的 $4$ 个苹果中，任取 $2$ 个，求重量在 $[80,85)$ 和 $[95,100)$ 中各有 $1$ 个的概率。

\Needspace{0.4\textheight}
\end{question}
\begin{question}
\gkti{18}如图①，在边长为 $1$ 的等边三角形 $ABC$ 中，$D$，$E$ 分别是 $AB$，$AC$ 上的点，$AD=AE$，$F$ 是 $BC$ 的中点，$AF$ 与 $DE$ 交于点 $G$。将 $\triangle ABF$ 沿 $AF$ 折起，得到如图②所示的三棱锥 $A-BCF$，其中 $BC=\dfrac{\sqrt2}2$。
\nopagebreak[4]
\begin{center}% 2013 广东卷文科第 18 题图①：等边三角形 ABC 中的折痕（AF 与 DE 交于 G）。
\begin{tikzpicture}[font=\small, >=Stealth, line width=0.4pt]
  \coordinate (B) at (0,0);
  \coordinate (C) at (3,0);
  \coordinate (A) at (1.5,2.5981);
  \coordinate (F) at (1.5,0);
  \coordinate (D) at (0.5,0.8660);
  \coordinate (E) at (2.5,0.8660);
  \coordinate (G) at (1.5,0.8660);
  \draw (A) -- (B) -- (C) -- cycle;
  \draw (A) -- (F);
  \draw (D) -- (E);
  \fill (G) circle (1.1pt);
  \node[above] at (A) {$A$};
  \node[below left] at (B) {$B$};
  \node[below right] at (C) {$C$};
  \node[below] at (F) {$F$};
  \node[left] at (D) {$D$};
  \node[right] at (E) {$E$};
  \node[above left] at (G) {$G$};
  \node (lbl) at (1.5,-0.85) {图};
  \node[circle, draw, inner sep=0.4pt, font=\scriptsize, right=1pt of lbl] {1};
\end{tikzpicture}
\end{center}
\nopagebreak[4]
\nopagebreak[4]
\begin{center}% 2013 广东卷文科第 18 题图②：折叠后的三棱锥 A-BCF（D、E、G 为折痕上的点）。
\begin{tikzpicture}[font=\small, >=Stealth, line width=0.4pt]
  \coordinate (A) at (1.55,3.60);
  \coordinate (B) at (0,0);
  \coordinate (C) at (3.40,0.55);
  \coordinate (F) at (1.45,0.75);
  \coordinate (D) at (0.5167,1.20);
  \coordinate (E) at (2.7833,1.5667);
  \coordinate (G) at (1.4833,1.70);
  % 被遮挡的棱与辅助线
  \draw[dashed] (A) -- (F);
  \draw[dashed] (F) -- (B);
  \draw[dashed] (F) -- (C);
  \draw[dashed] (G) -- (D);
  \draw[dashed] (G) -- (E);
  % 可见的棱
  \draw (A) -- (B) -- (C) -- (A);
  \draw (D) -- (E);
  \fill (F) circle (1.1pt);
  \fill (G) circle (1.1pt);
  \node[above] at (A) {$A$};
  \node[below left] at (B) {$B$};
  \node[right] at (C) {$C$};
  \node[below] at (F) {$F$};
  \node[left] at (D) {$D$};
  \node[right] at (E) {$E$};
  \node[above left] at (G) {$G$};
  \node (lbl) at (1.55,-0.75) {图};
  \node[circle, draw, inner sep=0.4pt, font=\scriptsize, right=1pt of lbl] {2};
\end{tikzpicture}
\end{center}
\nopagebreak[4]
（1） 证明：$DE\parallel$ 平面 $BCF$；

（2） 证明：$CF\perp$ 平面 $ABF$；

（3） 当 $AD=\dfrac23$ 时，求三棱锥 $F-DEG$ 的体积 $V_{F-DEG}$。

\end{question}
\begin{question}
\gkti{19}设各项均为正数的数列 $\{a_n\}$ 的前 $n$ 项和为 $S_n$，满足 $4S_n=a_{n+1}^2-4n-1$，$n\in\mathbf N^*$，且 $a_2$，$a_5$，$a_{14}$ 构成等比数列。

（1） 证明：$a_2=\sqrt{4a_1+5}$；

（2） 求数列 $\{a_n\}$ 的通项公式；

（3） 证明：对一切正整数 $n$，有 $\dfrac1{a_1a_2}+\dfrac1{a_2a_3}+\cdots+\dfrac1{a_na_{n+1}}<\dfrac12$。

\end{question}
\begin{question}
\gkti{20}已知抛物线 $C$ 的顶点为原点，其焦点 $F(0,c)$（$c>0$）到直线 $l:x-y-2=0$ 的距离为 $\dfrac{3\sqrt2}2$。设 $P$ 为直线 $l$ 上的点，过点 $P$ 作抛物线 $C$ 的两条切线 $PA$，$PB$，其中 $A$，$B$ 为切点。

（1） 求抛物线 $C$ 的方程；

（2） 当点 $P(x_0,y_0)$ 为直线 $l$ 上的定点时，求直线 $AB$ 的方程；

（3） 当点 $P$ 在直线 $l$ 上移动时，求 $|AF|\cdot|BF|$ 的最小值。

\end{question}
\begin{question}
\gkti{21}设函数 $f(x)=x^3-kx^2+x$（$x\in\mathbf R$）。

（1） 当 $k=1$ 时，求函数 $f(x)$ 的单调区间；

（2） 当 $k<0$ 时，求函数 $f(x)$ 在 $[k,-k]$ 上的最小值 $m$ 和最大值 $M$。

\end{question}